%----------------------------------------------------------------------------------------
%	PACKAGES AND OTHER DOCUMENT CONFIGURATIONS
%----------------------------------------------------------------------------------------

\documentclass[11pt, a4paper, twocolumn]{article}

%----------------------------------------------------------------------------------------
%	PACKAGES AND OTHER DOCUMENT CONFIGURATIONS
%----------------------------------------------------------------------------------------

\usepackage[english]{babel} % English language hyphenation

\usepackage{microtype} % Better typography

\usepackage{amsmath,amsfonts,amsthm} % Math packages for equations

\usepackage[svgnames]{xcolor} % Enabling colors by their 'svgnames'
\definecolor{TUBlue}{RGB}{0,48,94} %TU_Dresden Corporate coloring

\usepackage[hang, small, labelfont=bf, up, textfont=it]{caption} % Custom captions under/above tables and figures

\usepackage{booktabs} % Horizontal rules in tables

\usepackage{lastpage} % Used to determine the number of pages in the document (for "Page X of Total")

\usepackage{graphicx} % Required for adding images

\usepackage{enumitem} % Required for customising lists
\setlist{noitemsep} % Remove spacing between bullet/numbered list elements

\usepackage[export]{adjustbox} % Easy image ragging

\usepackage{sectsty} % Enables custom section titles
\allsectionsfont{\usefont{OT1}{phv}{b}{n}} % Change the font of all section commands (Helvetica)

%----------------------------------------------------------------------------------------
%	MARGINS AND SPACING
%----------------------------------------------------------------------------------------

\usepackage{geometry} % Required for adjusting page dimensions

\geometry{
	top=1cm, % Top margin
	bottom=1.5cm, % Bottom margin
	left=2cm, % Left margin
	right=2cm, % Right margin
	includehead, % Include space for a header
	includefoot, % Include space for a footer
	%showframe, % Show how the type block is set on the page
}

\setlength{\columnsep}{7mm} % Column separation width

%----------------------------------------------------------------------------------------
%	FONTS
%----------------------------------------------------------------------------------------

\usepackage[T1]{fontenc} % Output font encoding for international characters
\usepackage[utf8]{inputenc} % Required for inputting international characters

\usepackage{XCharter} % Use the XCharter font

%----------------------------------------------------------------------------------------
%	HEADERS AND FOOTERS
%----------------------------------------------------------------------------------------

\usepackage{fancyhdr} % Needed to define custom headers/footers
\pagestyle{fancy} % Enables the custom headers/footers

\renewcommand{\headrulewidth}{0.0pt} % No header rule
\renewcommand{\footrulewidth}{0.4pt} % Thin footer rule

\renewcommand{\sectionmark}[1]{\markboth{#1}{}} % Removes the section number from the header when \leftmark is used

%\nouppercase\leftmark % Add this to one of the lines below if you want a section title in the header/footer

% Headers
\lhead{} % Left header
\chead{\textit{\thetitle}} % Center header - currently printing the article title
\rhead{} % Right header

% Footers
\lfoot{} % Left footer
\cfoot{} % Center footer
\rfoot{\footnotesize Page \thepage\ of \pageref{LastPage}} % Right footer, "Page 1 of 2"

\fancypagestyle{firstpage}{ % Page style for the first page with the title
	\fancyhf{}
	\renewcommand{\footrulewidth}{0pt} % Suppress footer rule
}

%----------------------------------------------------------------------------------------
%	TITLE SECTION
%----------------------------------------------------------------------------------------

\newcommand{\authorstyle}[1]{{\large\usefont{OT1}{phv}{b}{n}\color{Black}#1}} % Authors style (Helvetica)

\newcommand{\institution}[1]{{\footnotesize\usefont{OT1}{phv}{m}{sl}\color{Black}#1}} % Institutions style (Helvetica)

\usepackage{titling} % Allows custom title configuration

\newcommand{\HorRule}{\color{TUBlue}\rule{\linewidth}{1.5pt}} % Defines the horizontal rule around the title

\pretitle{
	\vspace{-30pt} % Move the entire title section up
	\includegraphics[width=0.3\textwidth, right]{img/tud.pdf}
	\HorRule\vspace{10pt} % Horizontal rule before the title
	\fontsize{24}{28}\usefont{OT1}{phv}{b}{n}\selectfont % Helvetica
	\color{Black} % Text colour for the title and author(s)
}

\posttitle{\par\vskip 15pt} % Whitespace under the title

\preauthor{} % Anything that will appear before \author is printed

\postauthor{ % Anything that will appear after \author is printed
	\vspace{0pt} % Space before the rule
	\par\HorRule % Horizontal rule after the title
	\vspace{20pt} % Space after the title section
}

%----------------------------------------------------------------------------------------
%	ABSTRACT
%----------------------------------------------------------------------------------------

\usepackage{lettrine} % Package to accentuate the first letter of the text (lettrine)
\usepackage{fix-cm}	% Fixes the height of the lettrine

\newcommand{\initial}[1]{ % Defines the command and style for the lettrine
	\lettrine[lines=3,findent=4pt,nindent=0pt]{% Lettrine takes up 3 lines, the text to the right of it is indented 4pt and further indenting of lines 2+ is stopped
		\color{TUBlue}% Lettrine colour
		{#1}% The letter
	}{}%
}

\usepackage{xstring} % Required for string manipulation

\newcommand{\lettrineabstract}[1]{
	\StrLeft{#1}{1}[\firstletter] % Capture the first letter of the abstract for the lettrine
	\initial{\firstletter}\textbf{\StrGobbleLeft{#1}{1}} % Print the abstract with the first letter as a lettrine and the rest in bold
}

%----------------------------------------------------------------------------------------
%	BIBLIOGRAPHY
%----------------------------------------------------------------------------------------

\usepackage[backend=biber,style=authoryear,natbib=true]{biblatex} % Use the bibtex backend with the authoryear citation style (which resembles APA)

\addbibresource{literature.bib} % The filename of the bibliography

\usepackage[autostyle=true]{csquotes} % Required to generate language-dependent quotes in the bibliography
 % Specifies the document structure and loads requires packages

%----------------------------------------------------------------------------------------
%	ARTICLE INFORMATION
%----------------------------------------------------------------------------------------

\title{Forecasting with parametric models and neural networks: a comparison (exposé)} % The article title

\author{
	\authorstyle{Niklas Paulig\textsuperscript{1,2,3}} % Authors
	\newline\newline % Space before institutions
	\textsuperscript{1}\institution{Technical University of Dresden, Dresden, Germany}\\
	\textsuperscript{2}\institution{Mat.Nr: 4064033}\\
	\textsuperscript{3}\institution{Master VWL, Fak. Wirtschaftswissenschaften}
}


\date{\today} % Add a date here 
%----------------------------------------------------------------------------------------

\begin{document}

\maketitle % Print the title

\thispagestyle{firstpage} % Apply the page style for the first page (no headers and footers)

%----------------------------------------------------------------------------------------
%	ABSTRACT
%----------------------------------------------------------------------------------------

\lettrineabstract{For this master thesis, I choose 14 stock indices from around the world to compare the forecasting abilities of different parametric statistical models against the forecasting abilities of artificial neural networks (ANNs). The focus will lie on the inner workings and parameterizations of those two tools as well as accuracy and goodness-of-fit measures. The research question asks whether black-box approaches like neural networks can outperform profound statistical forecasting models.}

%----------------------------------------------------------------------------------------
%	ARTICLE CONTENTS
%----------------------------------------------------------------------------------------

\section{Idea}

When performing time-series forecasting the usual goal is not only to predict future values from a given set of data but also to get an understanding of the data generating process. Both forcasting branches I cover in this thesis do this, but very differently.

For classical time-series analysis, I will set up an exponential smoothing model based on the works of \citet{brown1962smoothing}, \citet{holt1957forecasting} and \citet{winters1960forecasting} that will act as a baseline. Various error metrics to measure out-of-sample performance will be provided, such as RMSE, MAE and MAPE. 

In a next step, the more advanced \citet{box1976time} methodology is used to fit an ARIMA model to the data and compare it to the naive exponential smoothing in terms of accuracy and robustness. As the 14 time-series are stock indices with fluctuating variances I will also specify and fit a generalized ARCH  model to the series in order to account for its volatility clustering (\citet{bollerslev1994arch} based on the work of \citet{engle1982autoregressive}).

The classic models will be compared using relative entropy estimators like the Akaike information criterion and other related estimators such as Hannan-Quinn or the Schwarz information criterion. Unfortunately due to the nature of artificial neural networks not using likelihood functions, those cannot be compared to the classical models in this way. 

The second part of the thesis will introduce artificial neural networks, that substantially differ from the classical models in the way they work and produce forecasts. The development of neural networks will be traced back to the Perceptron algorithm (\citet{rosenblatt1957perceptron}) and the idea of a neuron will be explained. 

As a benchmark for feed-forward neural networks a multi-layer perceptron (MLP) with different parametrizations is trained on the data and out-of-sample performance in comparison to classical models is discussed. In a final step, a more complex recurrent neural network (RNN) (see for example \citet{williams1989learning} or \citet{werbos1990backpropagation}), a Long Short-Term Memory (LSTM) model after \citet{hochreiter1997long} is fitted and compared in order to get a more holistic overview over the forecasting abilities of simple and more complex artificial neural networks.

\section{Structure}

The thesis structure is planned to contain three chapters. The first chapter introduces the different inference and forecast models and explains their origin, use cases and inner workings. I will thereby start with the classical models (exp. smoothing, ARIMA, GARCH) and continue with the neural networks. The last part of the first chapter will introduce the various error and performance metrics and explain their behavior. 

Chapter two introduces the data-set and firstly describes the pre-fit transformations and statistical tests and secondly the fitting process itself. The models are then compared to each other with the above-introduced out-of-sample performance metrics. 

The last chapter draws a conclusion to the question whether the rather complicated and opaque neural networks can outperform classical statistical analysis, and if they do, to what magnitude. 



%----------------------------------------------------------------------------------------
%	BIBLIOGRAPHY
%----------------------------------------------------------------------------------------

\printbibliography[title={References}] % Print the bibliography, section title in curly brackets

%----------------------------------------------------------------------------------------

\end{document}
